% page 512 of the facsimile (image p-511 of the scan)
% block 160.0 x 250.6 mm ; left margin 8.0 mm
\pgc{-12.700mm}{0.000mm}{
\l{\cel{2}504.}
\l{}
\l{\sou{rovro}. (\sou{bot.)}\cel{2}Querko blanka, min granda kam la querko or-}
\l{\cel{3}dinara. - L.\cel{1}\sou{quercus robur}. - FIS.}
\l{}
\l{\sou{royalismo}. Mento monarkiista.(\textquotedbl{}En rejio, royalismo esas o}
\l{\cel{3}partiso, od opiniono\textquotedbl{}.) - DEFR.}
\l{}
\l{\sou{rozo.}\cel{2}(\sou{bot.}) Flori di arbusto ek la familio \textquotedbl{}rozacei\textquotedbl{}, qua}
\l{\cel{3}odoras dolcege, e di qua la tipo primitiva esas ye koloro}
\l{\cel{3}delikate pale-reda. - L.\cel{1}\sou{rosa}. -\cel{1}\sou{Vento-rozo}\cel{1}:Plako qua}
\l{\cel{3}subtenas \textquotedbl{}rozego\textquotedbl{} 32-folia, qua dividas la cirklokurvo}
\l{\cel{3}ad 32 parti inter-egala, kun la indiko, en singla de ta}
\l{\cel{3}dividuri, di la rumbo. - DEFIRS.}
\l{}
\l{\sou{rozario}. Granda chapleto, ek dek-e-kin deki de grani, sur sin-}
\l{\cel{3}gla de qui onu recitas (che la katoliki) un \textquotedbl{}Ave\textquotedbl{} - singla}
\l{\cel{3}de qui pre-iresas da grano plu grosa sur qua onu recitas}
\l{\cel{3}un \textquotedbl{}Pater\textquotedbl{}. - DEFIS.}
\l{}
\l{\sou{rozeolo}. (\sou{patol.}) Pel-erupto febla, karakterizata da mikra}
\l{\cel{3}makuli rozea, kurta-dura, sen febro, qua eventas precipue}
\l{\cel{3}che la infanteti e la infanti. - EFIS.}
\l{}
\l{\sou{rozlauro}.\cel{2}(\sou{bot.}) Arboreto ek la familio \textquotedbl{}apocinei\textquotedbl{}. -}
\l{\cel{3}L.\cel{1}\sou{nerium oleander}.}
\l{}
\l{\sou{rubando}. - Mikra bendo streta de stofo fila, lana, silka,}
\l{\cel{3}velura, de celofano, de plastiko. - EF.}
\l{}
\l{\sou{rubarbo}. (\sou{bot.}) Planto ek la familio \textquotedbl{}poligonei\textquotedbl{}, kun larja}
\l{\cel{3}folii, di qua la radiki uzesas kom tonizivo e purgivo, e}
\l{\cel{3}la pedunkli, por facar kompoti e tarti. - L.\cel{1}\sou{rheum rhubarba}.}
\l{\cel{3}- DEFIS.}
\l{}
\l{\sou{rubio}. (\sou{bot.}) Planto ek la familio \textquotedbl{}rubimacei\textquotedbl{}, di qua la}
\l{\cel{3}radiko, kande sikigita e pulverigita, donas tintivo reda.}
\l{\cel{3}L.\cel{1}\sou{rubia tinctorum}. - ISL.}
\l{}
\l{\sou{rubidio}. (\sou{kemio)}. Metalo alkalioza, qua existas en ula aqui}
\l{\cel{3}minerala, en ula vejetanti (betravo, tabako, teo, kafeo),}
\l{\cel{3}en la lepidolito, bazalto, e c., e deskovrita samatempe}
\l{\cel{3}kam cezio, da Kirchhoff e Bunsen. -\cel{1}\sou{Simbolo kemiala}:Rh.}
\l{}
\l{\sou{rubino}. Lapido, reda bele e diafane. DEFIRS.}
\l{}
\l{\sou{rubriko}. En revuo, en jurnalo : titulo generala (kom exemplo:}
\l{\cel{3}\textquotedbl{}Progreso\textquotedbl{}) kun tituli specala (kom exemplo, ek \textquotedbl{}Progreso\textquotedbl{} -}
\l{\cel{3}rubrikin esas\textquotedbl{}Questioni linguala\textquotedbl{}, \textquotedbl{}Kroniko\textquotedbl{}, \textquotedbl{}Bibliogra-}
\l{\cel{3}fio\textquotedbl{}, \textquotedbl{}Diversaji\textquotedbl{}.) - DEFIS.}
\l{}
\l{\sou{ruda}. Di qua la acepto-maniero esas ne-delikata, kelke male-}
\l{\cel{3}traktanta, mem se nur parole. - DEFI.}
\l{}
\l{rudimento.\cel{2}I. La komenc-elementi di arto, di cienco. - II.}
\l{\cel{3}Mikra libro qua kontenas la komenc-elementi di la grama-}
\l{\cel{3}tiko Latina. - III. Komenco-lineamento di la strukturo}
\l{\cel{3}di organo. - DEFIRS.}
}
